\documentclass[12pt]{article}
\usepackage{amsmath,amsfonts,fullpage,amsthm,graphicx}

\newtheorem*{theorem}{Theorem}
\newtheorem*{fact}{Fact}
\newtheorem*{goal}{Goal}
\newtheorem*{idea}{Idea}
\newtheorem*{defn}{Definition}
\newtheorem*{ex}{Example}

%Style section
%\setlength{\textheight}{10in}
%\setlength{\textwidth}{7in}
%\hoffset=-0.75in
 \pagestyle{plain}
% \setlength{\topmargin}{0in}
% \setlength{\headsep}{0in}
% \setlength{\oddsidemargin}{0in}
% \setlength{\evensidemargin}{0in}

\begin{document}


\pagestyle{empty} %use this to get rid of page numbers

\noindent
{Math 166  \hskip 1.5 truein  Name:\ \hrulefill}

\noindent
%\vskip 0.3truein
%\bigskip
%\noindent
{A study of series}

\noindent  \hskip 3.5 truein  Section Number:\ \hrulefill
\medskip

\begin{goal} \rm
Explore even more methods of determining whether series converge or diverge. 
\end{goal}
%\smallskip

\begin{enumerate}
%\setcounter{enumi}{-1}
\item Name all the tests for convergence/divergence of series in our series toolbox. Put a * by the tests that only work for series with positive terms.
%\item %A \textbf{sequence} is a~ \rule{.8in}{.4pt} ~of numbers.  A sequence converges if its~ \rule{.8in}{.4pt} ~have a limit.
%\item A \textbf{series} is a~ \rule{.6in}{.4pt} ~of numbers. A series converges if its~ \rule{.8in}{.4pt} \rule{.2in}{.4pt} \rule{.8in}{.4pt} %\rule{.5in}{.4pt} ~has a limit.
%\item
%\end{enumerate}
\vfill

\item Use any of the tests you listed in 1.\ to determine whether the following series converge or diverge. Be sure to verify the hypotheses, state your conclusion, and say which test you used. (Feel free to approach to these problems together with your group, each trying different tests and sharing with each other what worked.)
\begin{enumerate}
\item $\displaystyle\sum_{n=1}^{\infty} (-1)^n e^{-3n}$

\vfill
\newpage

%\item $\displaystyle\sum_{n=1}^{\infty} \frac{e^n}{n!}$
%
%\vfill
\hrule\ \\
{\bf Theorem.} If a series \emph{converges absolutely}, then it converges.  That is, if the series $\sum_{n = 0}^\infty |a_n|$ converges, then so does the series $\sum_{n = 0}^\infty a_n$.\\
\hrule
\item $\displaystyle\sum_{n=1}^{\infty} \frac{\cos(n)}{n^2}$ (Hint: Use the theorem above.)

\vfill

%\item $\displaystyle\sum_{n=1}^{\infty} \frac{5^n}{n^n}$ %(Hint: When you see something to the $n$th power in a series, it's often a good idea to use the root test.)
%
%\vfill


\item $\displaystyle\sum_{n=1}^{\infty} (-1)^{n-1} \frac{7^n+1}{2^n}$

\vfill

\newpage

\item $\displaystyle\sum_{n=1}^{\infty} (-1)^n \ln(n)$

\vfill
%\item $\displaystyle\sum_{n=1}^{\infty} \frac{e^n}{n!}$

\end{enumerate} 

\item Approximate $\displaystyle\sum_{n=1}^{\infty} (-1)^n \displaystyle\frac{1}{n^3}$
to two decimal places.  Hint: for alternating series, the error of a partial sum is bounded by the next term in the sequence.
\vfill

\pagebreak
\item Do the same as problem 2 for the following series. %$\displaystyle\sum_{n=1}^\infty \cos(\pi n) \frac{\mathrm{ln}(n)}{n}$
\begin{itemize}
\item $\displaystyle\sum_{n=1}^\infty \cos(\pi n) \frac{\mathrm{ln}(n)}{n}$
\vfill

\item $\displaystyle\sum_{n=1}^{\infty} \frac{n^3-5n+1}{7n^5+2}$

\vfill


\item $\displaystyle\sum_{k=2}^{\infty} \sqrt{\frac{k}{k^3+1}}$

% $\displaystyle\sum_{n=1}^\infty (-2)^n \frac{n}{(n+2)2^n}$
%\vfill

%\item $\displaystyle\sum_{n=1}^\infty \frac{(-1)^2}{-(n+2)^2}$
\end{itemize}
\vfill
\end{enumerate}

\end{document}
